%% ========================================================================
%% CAPÍTULO 1 - INTRODUÇÃO
%% ========================================================================

\capitulo{intro}{Introdução}

A mobilidade urbana constitui um dos principais desafios das cidades contemporâneas, especialmente em regiões metropolitanas de países em desenvolvimento. O crescimento populacional acelerado, associado à expansão urbana desordenada, tem gerado problemas complexos relacionados ao transporte de pessoas e mercadorias \cite{vasconcellos2012}.

No Brasil, as questões de mobilidade urbana ganharam destaque nas últimas décadas devido ao aumento significativo da frota de veículos particulares e à inadequação dos sistemas de transporte público. Segundo dados do \sigla{IBGE} \cite{ibge2022}, a frota nacional de veículos cresceu 300\% entre 1990 e 2020, enquanto a população urbana aumentou apenas 60\% no mesmo período.

\section{Contextualização}

A Região Metropolitana de Brasília, composta pelo Distrito Federal e municípios do entorno, apresenta características peculiares que influenciam diretamente os padrões de mobilidade urbana. O planejamento original da capital, concebido por Lucio Costa na década de 1950, priorizou o transporte individual motorizado, resultando em uma estrutura urbana dispersa e dependente do automóvel \cite{costa2018}.

Atualmente, a região enfrenta desafios significativos relacionados à mobilidade, incluindo:

\begin{itemize}
    \item Congestionamentos frequentes nos principais corredores viários;
    \item Baixa qualidade do transporte público;
    \item Infraestrutura inadequada para pedestres e ciclistas;
    \item Altos índices de poluição atmosférica;
    \item Segregação socioespacial relacionada à acessibilidade.
\end{itemize}

\section{Justificativa}

A necessidade de estudos sobre mobilidade urbana sustentável na Região Metropolitana de Brasília justifica-se por diversos fatores. Primeiro, a região apresenta um dos maiores índices de motorização do país, com aproximadamente 0,7 veículos por habitante \cite{detran2023}. Segundo, os custos sociais e ambientais do modelo atual de mobilidade são crescentes, incluindo tempo perdido em congestionamentos, acidentes de trânsito e emissões de gases poluentes.

Além disso, a implementação de políticas públicas de mobilidade urbana requer fundamentação técnica e científica adequada. A Lei Federal nº 12.587/2012, que institui a Política Nacional de Mobilidade Urbana, estabelece diretrizes para o desenvolvimento de sistemas de transporte mais sustentáveis e equitativos \cite{brasil2012}.

\section{Problema de Pesquisa}

O problema central desta pesquisa pode ser formulado da seguinte forma: \textbf{Como promover a transição para um sistema de mobilidade urbana mais sustentável na Região Metropolitana de Brasília, considerando as características específicas da região e as demandas da população?}

Esta questão principal desdobra-se em problemas específicos:

\begin{enumerate}
    \item Quais são as características atuais dos padrões de mobilidade na região?
    \item Quais os principais obstáculos para a adoção de modos de transporte mais sustentáveis?
    \item Que estratégias podem ser implementadas para melhorar a sustentabilidade do sistema de transporte?
    \item Como integrar os diferentes modais de transporte de forma eficiente?
\end{enumerate}

\section{Objetivos}

\subsection{Objetivo Geral}

Analisar as condições atuais de mobilidade urbana na Região Metropolitana de Brasília e propor estratégias integradas para o desenvolvimento de um sistema de transporte mais sustentável e eficiente.

\subsection{Objetivos Específicos}

\begin{enumerate}
    \item Caracterizar os padrões atuais de mobilidade urbana na região, identificando os principais modais utilizados e suas características operacionais;
    
    \item Avaliar a qualidade dos serviços de transporte público, considerando aspectos como acessibilidade, frequência, conforto e integração;
    
    \item Analisar os impactos ambientais, sociais e econômicos dos diferentes modais de transporte;
    
    \item Identificar as principais barreiras para a adoção de modos de transporte mais sustentáveis;
    
    \item Propor um conjunto de estratégias integradas para promover a mobilidade urbana sustentável;
    
    \item Desenvolver indicadores para monitoramento e avaliação das políticas de mobilidade urbana.
\end{enumerate}

\section{Hipóteses}

As hipóteses que orientam esta pesquisa são:

\begin{enumerate}
    \item A predominância do transporte individual motorizado na Região Metropolitana de Brasília está relacionada às deficiências do sistema de transporte público e à inadequação da infraestrutura para modos não motorizados;
    
    \item A implementação de estratégias integradas de mobilidade urbana pode reduzir significativamente a dependência do automóvel e melhorar a sustentabilidade do sistema de transporte;
    
    \item A integração física, tarifária e operacional entre os diferentes modais de transporte é fundamental para o sucesso de políticas de mobilidade sustentável.
\end{enumerate}

\section{Estrutura do Trabalho}

Este trabalho está estruturado em quatro capítulos principais:

O \refCap{intro} apresenta a contextualização do problema, justificativa, objetivos e hipóteses da pesquisa.

O \refCap{metodologia} descreve os procedimentos metodológicos adotados, incluindo a caracterização da área de estudo, métodos de coleta e análise de dados.

O \refCap{resultados} apresenta e discute os resultados obtidos, incluindo a caracterização dos padrões de mobilidade, avaliação dos sistemas de transporte e proposição de estratégias.

O \refCap{conclusao} sintetiza as principais conclusões, contribuições da pesquisa e recomendações para trabalhos futuros.

\section{Delimitação do Estudo}

Esta pesquisa tem como recorte espacial a Região Metropolitana de Brasília, incluindo o Distrito Federal e os municípios do entorno que fazem parte da \sigla{RIDE}. O recorte temporal abrange o período de 2020 a 2023, permitindo uma análise das transformações recentes nos padrões de mobilidade.

O estudo foca nos deslocamentos de pessoas, não abordando especificamente o transporte de cargas. Embora reconheça-se a importância da logística urbana para a mobilidade sustentável, esta temática não será aprofundada devido às limitações de escopo da pesquisa.